% Einheit 3 · Kreisteile
\setcounter{aufgabe}{19}
\einheitenkopf*{\hypertarget{e3}{}Einheit 3 von 3 · Kreisteile}
\zweigzeile{Hier lernst du, mit Kreisausschnitten zu rechnen: Anteil aus dem Winkel, Bogenlänge, Fläche des Ausschnitts und Kreisring · neu oder schon bekannt – je nach Buch · P10 · baut auf: Kreisumfang (Einheit 1), Kreisfläche (Einheit 2), mit dem Vollwinkel rechnen, Anteile als Bruch und Prozent}

\begin{aufgabe}{Ich erkenne, welcher Teil des Kreises grau ist. Kreuze den passenden Anteil an. Rechne nicht.}
\begin{minipage}[t]{3.1cm}
a)\\[1mm]
\begin{kreis}[1.1]\mittelpunkt{}\sektor{90}{270}{}\end{kreis}\\[1mm]
\kreuz{$\dfrac12$}\\[1.5mm] \kreuz{$\dfrac13$}\\[1.5mm] \kreuz{$\dfrac14$}
\end{minipage}%
\begin{minipage}[t]{3.1cm}
b)\\[1mm]
\begin{kreis}[1.1]\mittelpunkt{}\sektor{0}{90}{}\end{kreis}\\[1mm]
\kreuz{$\dfrac12$}\\[1.5mm] \kreuz{$\dfrac14$}\\[1.5mm] \kreuz{$\dfrac18$}
\end{minipage}%
\begin{minipage}[t]{3.1cm}
c)\\[1mm]
\begin{kreis}[1.1]\mittelpunkt{}\sektor{90}{210}{}\end{kreis}\\[1mm]
\kreuz{$\dfrac13$}\\[1.5mm] \kreuz{$\dfrac14$}\\[1.5mm] \kreuz{$\dfrac16$}
\end{minipage}%
\begin{minipage}[t]{3.1cm}
d)\\[1mm]
\begin{kreis}[1.1]\mittelpunkt{}\sektor{0}{270}{}\end{kreis}\\[1mm]
\kreuz{$\dfrac23$}\\[1.5mm] \kreuz{$\dfrac34$}\\[1.5mm] \kreuz{$\dfrac45$}
\end{minipage}%
\begin{minipage}[t]{3.1cm}
e)\\[1mm]
\begin{kreis}[1.1]\mittelpunkt{}\sektor{30}{90}{}\end{kreis}\\[1mm]
\kreuz{$\dfrac13$}\\[1.5mm] \kreuz{$\dfrac16$}\\[1.5mm] \kreuz{$\dfrac18$}
\end{minipage}
\end{aufgabe}

\begin{aufgabe}{Ich kann den Anteil eines Kreisausschnitts aus dem Mittelpunktswinkel berechnen. Der Anteil ist der Mittelpunktswinkel $\alpha$ geteilt durch 360°.}
\beispiel{\alpha = 30° \quad\rightarrow\quad \text{Anteil} &= 30 : 360 = \tfrac{1}{12}}
Gib den Anteil als gekürzten Bruch an.
\begin{teilezwei}
\tz $\alpha = 90°$ \quad Anteil: \leerfeld & \tz $\alpha = 180°$ \quad Anteil: \leerfeld \\
\tz $\alpha = 60°$ \quad Anteil: \leerfeld & \tz $\alpha = 45°$ \quad Anteil: \leerfeld \\
\end{teilezwei}
Gib den Anteil in Prozent an. Runde, wenn nötig, auf eine Stelle nach dem Komma.
\begin{teilezwei}
\tz $\alpha = 36°$ \quad Anteil: \leerfeld[\%] & \tz $\alpha = 50°$ \quad Anteil: \leerfeld[\%] \\
\end{teilezwei}
\end{aufgabe}

\begin{aufgabe}{Ich kann aus dem Anteil den Mittelpunktswinkel berechnen. Berechne den Mittelpunktswinkel $\alpha$.}
\beispiel{\text{Anteil } \tfrac{1}{12} \quad\rightarrow\quad \alpha &= 360° : 12 = 30°}
\begin{teilezwei}
\tz Anteil $\frac18$ \quad \feld[°]{\alpha} & \tz Anteil $\frac38$ \quad \feld[°]{\alpha} \\
\tz Anteil 30\,\% \quad \feld[°]{\alpha} & \\
\end{teilezwei}
\end{aufgabe}

\begin{aufgabe}{Ich kann Bogenlänge, Flächeninhalt und Umfang eines Kreisausschnitts berechnen. Runde auf eine Stelle nach dem Komma.}
\begin{kreis}[1.4]\mittelpunkt{M}\sektor{20}{100}{\alpha}\bogen{20}{100}{b}\radius{20}{r}\end{kreis}

\beispiel{r = 6\,\text{cm},\ \alpha = 60°: \quad b &= \tfrac{60}{360} \cdot 2 \cdot \pi \cdot 6\,\text{cm} \approx 6{,}3\,\text{cm} \\ A &= \tfrac{60}{360} \cdot \pi \cdot 6^2\,\text{cm}^2 \approx 18{,}8\,\text{cm}^2}
Berechne die Bogenlänge $b$.
\begin{teilezwei}
\tz $r = 7$\,cm, $\alpha = 90°$ \quad \feld[cm]{b} & \tz $r = 7$\,cm, $\alpha = 50°$ \quad \feld[cm]{b} \\
\end{teilezwei}
Berechne den Flächeninhalt $A$ des Ausschnitts.
\begin{teilezwei}
\tz $r = 7$\,cm, $\alpha = 50°$ \quad \feld[cm²]{A} & \\
\end{teilezwei}
Berechne den Umfang des Ausschnitts: Bogen und zwei Radien.
\begin{teilezwei}
\tz $r = 7$\,cm, $\alpha = 50°$ \quad \leerfeld[cm] & \\
\end{teilezwei}
Berechne Bogenlänge und Flächeninhalt.
\begin{teile}
\teil $r = 7{,}5$\,m, $\alpha = 132°$ \quad \feld[m]{b} \quad \feld[m²]{A}
\end{teile}
\end{aufgabe}

\begin{aufgabe}{Ich kann den Flächeninhalt eines Kreisrings berechnen. Ein Kreisring liegt zwischen einem großen Kreis mit dem Radius $r_a$ und einem kleinen Kreis mit dem Radius $r_i$ um denselben Mittelpunkt. Runde auf eine Stelle nach dem Komma.}
\begin{kreis}[1.8]\draw[fill=black!15] (0,0) circle (1.8); \draw[fill=white] (0,0) circle (1.0); \mittelpunkt{M}\radius{30}{r_a}\draw (0,0) -- (240:1.0) node[midway,left]{$r_i$};\end{kreis}

\beispiel{r_a = 9\,\text{cm},\ r_i = 6\,\text{cm}: \quad A &= \pi \cdot 9^2 - \pi \cdot 6^2 \approx 141{,}4\,\text{cm}^2}
\begin{teilezwei}
\tz $r_a = 8$\,cm, $r_i = 6$\,cm \quad \feld[cm²]{A} & \tz $r_a = 11$\,m, $r_i = 1$\,m \quad \feld[m²]{A} \\
\end{teilezwei}
\begin{teile}
\teil Der große Kreis hat den Durchmesser 22\,cm, der kleine den Durchmesser 16\,cm. \quad \feld[cm²]{A}
\teil Das Kreisdiagramm zeigt, welcher Teil einer Klasse mit dem Bus zur Schule kommt: der graue Teil mit 125°. Welcher Anteil der Klasse kommt nicht mit dem Bus? Gib den Anteil in Prozent an und runde auf eine Stelle. (P10 2025)
\end{teile}

\kreissektor{125}{$125°$}
\end{aufgabe}

\begin{aufgabe}{Ich finde den Fehler beim Anteil eines Kreisausschnitts. Ben soll angeben, welcher Anteil des Kreises zu einem Ausschnitt mit $\alpha = 150°$ gehört. Er rechnet so:}
\rechnung{\text{Anteil} &= 150 : 180 \\ &\approx 0{,}833 \\ &= 83{,}3\,\%}
\begin{teile}
\teil Finde den Fehler, benenne ihn und korrigiere die Rechnung.
\teil Berechne den Anteil für $\alpha = 162°$ in Prozent.
\end{teile}
\end{aufgabe}

\begin{aufgabe}{Ich kann begründen, welcher Teil eines Kreises zu einem Mittelpunktswinkel gehört. Begründe ohne Rechnung.}
\begin{teile}
\teil Warum gehört zu einem Halbkreis der Mittelpunktswinkel 180°?
\teil Lina sagt: „Wenn ich bei gleichem Radius den Mittelpunktswinkel verdopple, verdoppelt sich die Fläche des Ausschnitts.“ Stimmt das?
\teil Ole sagt: „Ein Ausschnitt mit 90° ist immer größer als ein Ausschnitt mit 60°.“ Stimmt das?
\end{teile}
\end{aufgabe}

\begin{aufgabe}{Ich kann mit einem Kreisausschnitt eine Fläche planen. Ein Rasensprenger spritzt 9\,m weit und dreht sich dabei hin und her. Er ist auf einen Winkel von 90° eingestellt.}
\begin{teile}
\teil Wie groß ist die Fläche, die er bewässert? Runde auf eine Stelle.
\teil Ein Rasenstück von 130\,m² soll bewässert werden. Auf welchen Winkel muss man den Sprenger ungefähr einstellen?
\teil Der Sprenger schafft höchstens 180°. Reicht er für das Rasenstück? Begründe.
\end{teile}
\end{aufgabe}
